\section{Sistema temporal dodecafásico orientado del TPK}
\label{sec:ciclo-dodecafasico-tpk-rev7}
\index{TPK!sistema dodecafásico}\index{periodicidad de orden 108}

\HMTClaim{HMT-R-TPK-GAMMA108}

\subsection{Operador de selección de ruta y transporte al observable dimensional}

La función ternaria de fase operativa es
\[
M_{\rm ph}(r)=
\begin{cases}
 +1,&r\in\{1,4,7\},\\
 -1,&r\in\{2,5,8\},\\
 0,&r\in\{3,6,9\}.
\end{cases}
\]
Sus tres valores son las entradas escalares con las que
\(\operatorname{Sel}_t\) activa, respectivamente, el cursor aditivo, el
cursor multiplicativo o el contador neutral. El muestreo estroboscópico es
\[
 \operatorname{Obs}_3(a)=(a_3,a_6,a_9,\ldots).
\]
La selección \(\operatorname{Sel}_t[\,M_{\rm ph}(g(t))\,]\) decide qué
evaluación opera; el transporte cardinal desplaza el estado; y
\(\operatorname{Obs}_3\) selecciona qué instantes se publican.

El transporte cardinal actúa por cuatro traslaciones
\[
 T_d(i,j)=(i,j)+\nu_d\pmod9,
\]
\[
 \nu_N=(-1,0),\qquad \nu_E=(0,1),\qquad
 \nu_S=(1,0),\qquad \nu_O=(0,-1).
\]
Los cursores aditivo y multiplicativo conservan posiciones y orientaciones
separadas. La coordenada TRIT, la función \(M_{\rm ph}\), el operador
\(\operatorname{Sel}_t\), las cuatro traslaciones y
\(\operatorname{Obs}_3\) son, por ello, objetos de tipos distintos.

\subsection{Calendario y registro dodecafásicos}

El núcleo topológico de fase no termina en la actualización de una posición
sobre \(\mathbb Z_9^2\). Su estado temporal mínimo conserva el orden de los
eventos, la orientación, el retorno y la memoria de las transiciones. Para
una órbita de \(108\) pasos se introduce la partición ordenada
\[
 E_m=\{9(m-1)+1,\ldots,9m\},
 \qquad m=1,\ldots,12,
\]
y se denota su soporte ordenado por
\[
 \Gamma_{108}=E_1\sqcup\cdots\sqcup E_{12},
 \qquad 108=12\cdot9.
\]
Equivalentemente,
\[
 \Theta_{108}=\mathbb Z/108\mathbb Z,\qquad
 \beta(\theta)=
 \left[\left\lfloor\frac{\widetilde\theta}{9}\right\rfloor\right]_{12},
 \qquad
 r(\theta)=1+(\theta\bmod9).
\]
La coordenada \(\beta\) selecciona uno de doce sectores y \(r\) localiza la
fase dentro de su ventana nonádica.

\begin{proposition}[Extensión central no escindida del reloj observable]
\label{prop:extension-central-108-no-escindida-rev8}
Escribiendo \(C_n=\mathbb Z/n\mathbb Z\), el calendario observable forma la
sucesión exacta
\[
 0\longrightarrow C_{12}
 \xrightarrow{\ \iota\ }C_{108}
 \xrightarrow{\ p\ }C_9
 \longrightarrow0,
 \qquad
 \iota([b]_{12})=[9b]_{108},\qquad
 p([\theta]_{108})=[\theta]_9.
\]
Para la sección conjuntista \(s([r]_9)=[r]_{108}\), con
\(0\leq r<9\), la costura viene dada por el cociclo de acarreo
\[
 c(r,r')=
 \left[\left\lfloor\frac{r+r'}{9}\right\rfloor\right]_{12},
 \qquad
 s(r)+s(r')=s([r+r']_9)+\iota\bigl(c(r,r')\bigr).
\]
La extensión no se escinde. Para la acción trivial,
\[
 H^2(C_9,C_{12})
 \simeq C_{12}/9C_{12}
 \simeq C_3,
\]
y \(c\) representa la clase generadora \([1]\).
\end{proposition}

\begin{proof}
La imagen de \(\iota\) es el subgrupo de los múltiplos de nueve en
\(C_{108}\), que coincide con \(\ker p\); de aquí la exactitud. La identidad
del cociclo es la división euclídea de \(r+r'\) por nueve. Si la extensión
central se escindiera, el grupo medio sería isomorfo a
\(C_{12}\times C_9\). Sin embargo,
\[
 \exp(C_{108})=108,\qquad
 \exp(C_{12}\times C_9)=\operatorname{mcm}(12,9)=36,
\]
lo que lo impide. La identificación cohomológica para un grupo cíclico con
acción trivial da \(H^2(C_9,C_{12})\simeq C_{12}/9C_{12}\), y el acarreo
unitario proyecta al generador de ese cociente.
\end{proof}

La no escisión tipa el retorno: nueve iteraciones cierran la coordenada en
\(C_9\) y trasladan una unidad en \(C_{12}\); ciento ocho iteraciones cierran ambas
coordenadas de este reloj. El calendario \(C_{108}\) es la proyección
temporal del estado enriquecido \(\mathcal X_{\rm TPK}^{\rm enr}\):
orientación, hoja, frontera, cilindro, genealogía y memoria profunda
permanecen en este último y pueden
distinguir historias con el mismo calendario observable.

Si \(g_0(\theta)=\theta\bmod9\), entonces
\[
 g_0(\theta+9)=g_0(\theta),\qquad
 \beta(\theta+9)=\beta(\theta)+1\pmod{12}.
\]
El testigo
\[
 (g_0,r,\beta,30^\circ\beta)(49)=(4,5,5,150^\circ),
\]
\[
 (g_0,r,\beta,30^\circ\beta)(58)=(4,5,6,180^\circ)
\]
muestra exactamente qué retorna y qué conserva memoria: la fase nonádica y
la posición interna coinciden, mientras el sector dodecafásico avanza uno.
El calibre regional \(\operatorname{diag}_c\) de las regiones de \(\pi\) no
es esta coordenada \(\beta\). Tampoco existe aquí un morfismo
\(30^\circ\beta\to h\in\langle90,120\rangle\); en particular,
\(150\notin\langle90,120\rangle\). El par \(49/58\) es un testigo de acarreo,
no un selector regional ni espectral.

Para una historia enriquecida \(x\) de TPK, sean \(\tau_m\) la subruta
ordenada sobre \(E_m\), \(\sigma_m\) su orientación, \(\kappa_m\) el acarreo
que cruza la frontera del bloque y \(\Lambda_m\) el estado acumulado del
registro. El objeto
\[
 \mathcal R_{12}
 =\bigl((E_m,\tau_m,\sigma_m,\kappa_m,\Lambda_m)\bigr)_{m=1}^{12}
\]
es el sistema temporal dodecafásico orientado de TPK. No constituye una cuarta
teoría añadida a APP, TRIT y TPK: es la forma temporal interna con la que
TPK ordena una órbita, conserva sus retornos y publica un registro
reversible. Se distinguen las proyecciones canónicas
\[
 \mathcal X_{\rm TPK}^{\rm enr}
 \longrightarrow \mathcal R_{12}
 \longrightarrow \Theta_{108}.
\]
La primera conserva \((\tau_m,\sigma_m,\kappa_m,\Lambda_m)\); la segunda
publica sólo la sucesión observable y no puede reconstruir por sí sola el
estado firmado. El registro enriquecido se denota \(\mathcal R_{12}\) y la
rotación de sectores se denota \(C_{12}\).

El cierre observable del calendario satisface
\[
 F_{\rm obs}^{108}=I.
\]
Las posiciones y orientaciones ya se restauran tras \(54\) iteraciones,
mientras \(\mathcal R_{12}\) distingue la posición temporal dentro de los
doce sectores. Aquí \(F_{\rm obs}\) es la permutación observable, mientras
\(\mathcal G_9\) es la poda dentro de la relación de extensión y frontera.
La monodromía es la propiedad conjunta
\[
 g(K+9)=g(K),\qquad v_{K+9}\ne v_K,
\]
no el nombre aislado de la poda. El cierre cronológico devuelve el calendario
observable; el retorno nonádico conserva la fase y modifica prefijo, cilindro
y memoria.

También se distinguen dos palabras de doce componentes:
\[
 U_{12}^{\rm cat}=U_6\Vert U_6,
 \qquad
 U_{12}^{\rm sgn}=\mathcal H_{12}(K).
\]
La primera concatena dos emisiones de longitud seis. La segunda es una
coordenada firmada del estado enriquecido y queda ligada reversiblemente a
\(K\) por la transformación dodecafásica \(\mathcal H_{12}\). Ni sus dominios
ni sus funciones son intercambiables; en particular, la palabra concatenada
no construye por sí sola la coordenada firmada.

\begin{definition}[Registro integral de los doce eventos]
\label{def:registro-dodecafasico-rev7}
Para una ruta de cifras \(d_1,\ldots,d_{108}\), sea
\[
 \mathcal D_{\rm evt}=\{2,4,6,8\},\qquad
 \Sigma_{12}=(+,+,+,-,-,-,+,+,+,-,-,-).
\]
Su palabra orientada de eventos es
\[
 e_m=\Sigma_{12}(m)\,
 \#\{t\in E_{m+1}:d_t\in\mathcal D_{\rm evt}\},
 \qquad 0\le m<12.
\]
Las dos semicoronas se emparejan mediante
\[
 p_i=e_i+e_{i+6},\qquad
 a_i=e_i-e_{i+6},\qquad 0\le i<6.
\]
Se definen el conteo bruto, las cinco diferencias respecto del sexto par y
las seis antisimetrías por
\[
 s_C=\sum_{i=0}^{5}p_i,\qquad
 M_i=p_i-p_5\quad(0\le i<5),\qquad
 A_6=(a_0,\ldots,a_5).
\]
El registro integral es
\[
 \mathscr L_{12}(e)
 =\bigl(s_C,M_0,\ldots,M_4,a_0,\ldots,a_5\bigr).
\]
\end{definition}

\begin{theorem}[Coordenadas exactas de la memoria dodecafásica]
\label{thm:memoria-1-5-6}
El mapa \(e\mapsto\mathscr L_{12}(e)\) es inyectivo sobre su imagen. La
inversa viene dada por
\[
 p_5=\frac{s_C-\sum_{i=0}^{4}M_i}{6},\qquad
 p_i=p_5+M_i\quad(0\le i<5),
\]
\[
 e_i=\frac{p_i+a_i}{2},\qquad
 e_{i+6}=\frac{p_i-a_i}{2}.
\]
En particular, la distribución de coordenadas es
\[
 \boxed{1+5+6=12}.
\]
\end{theorem}

\begin{proof}
Por definición,
\[
 s_C=\sum_{i=0}^{4}(p_5+M_i)+p_5
 =6p_5+\sum_{i=0}^{4}M_i.
\]
Esto determina \(p_5\) y después todos los \(p_i\). Sumar y restar las
ecuaciones \(p_i=e_i+e_{i+6}\) y \(a_i=e_i-e_{i+6}\) recupera los doce
eventos. La divisibilidad por seis y las seis condiciones de paridad
caracterizan la imagen integral; no intervienen una masa ni una unidad.
\end{proof}

La división por seis no se aplica durante el conteo y no es un coeficiente
primitivo. Primero se forma el conteo orientado bruto \(s_C\); después se
reconstruye el registro de doce eventos; sólo entonces aparece
\(s_C/6\) en el carácter. Equivalentemente,
\[
 W_\pm=6n_A\pm s_C
\]
son enteros y la factorización reticular posee forma normal de Smith
\((1,12)\). Esta precedencia será obligatoria en la ley de masas.

Los enteros \(90\) y \(120\) reaparecen más adelante en seis construcciones
de tipos distintos. Para impedir que el nombre del cardinal suplante al
objeto, se fija desde ahora el siguiente convenio:

\begin{enumerate}
 \item \(\mathcal E_{\rm dod}=(D_3,D_4,Q):\mathbb Z^{12}\to\mathbb Z^{25}\),
       extractor integral que reconstruye \(K\);
 \item \(D_{\rm raw}^{90,120}(q)\), diferencia diagnóstica de dos series en
       una carta fijada;
 \item \(\iota_{6,8}:\mathcal F_6\hookrightarrow\mathcal F_8\), inclusión
       combinatoria de cardinalidades \(90\) y \(120\);
 \item \(K_{6\to8}=12\Delta S_{90}-\Delta S_{120}\), funcional angular
       definido después de fijar canales, series y normalización;
 \item \(\mathfrak S=\langle90,120\rangle\), semigrupo infinito de alturas
       espectrales;
 \item \(\varepsilon_{\rm res}^{\circ}\), residuo de una identidad angular
       exacta, no una extracción finita de cualquiera de los cinco objetos
       anteriores.
\end{enumerate}

La coincidencia parcial de índices no establece igualdad de operadores,
dominios, codominios ni fuerza probatoria. En particular,
\(K_{6\to8}\) usa dos subcolas de la jerarquía, pero no es la jerarquía
completa; y ninguna de esas subcolas forma un bloque independiente de la ley
de masas.

\begin{remark}[Separación canónica de los objetos temporales]
La notación distingue el registro orientado \(\mathcal R_{12}\), el
calendario observable \(\Theta_{108}\), la rotación sectorial \(C_{12}\) y
la órbita de transporte \(\Gamma_{108}\), cada uno con su dominio, su
codominio y su ley de composición.
\end{remark}
